% A shifted main Pell parameter shares the existing B-4 mask register.
\section{A base-four congruence and the main Pell parameter}
\label{sec:main-base-four}

The third digit mask already computes $B-4$. This permits a useful joint
change: replace the base-two exponent test by a base-four test, and store
the main Pell parameter minus four. The two expressions $A-4$ and $A-B$
then share available values, eliminating one subtraction. The altered
Pell approximation lies strictly above its binomial expression; this
also makes the final rounding argument independent of parity.

Use the 107-operation system of Theorem~\ref{thm:best107}, with the
fixed modular-Sidon encoding and $L=2^{32}$ from
Section~\ref{sec:modular-support}. In this section $Y_P$ denotes
the Pell quantity $s n^2$, to distinguish it from the packed code.
Write
\[
\begin{gathered}
 N=n,\quad R=r,\quad U=wN^2,\quad Y_P=sN^2,\quad M=RY_P,\\
 a_0=M(U+1),\quad A=a_0+4,\quad Q=UM^2,\quad P=2Q+1,\\
 J=2R+1,\qquad W=bw,\qquad \xi=\frac{(U+1)^{2R}}{U^R}.
\end{gathered}
\]
The supplied positive unknown named $a$ in the certificate now means
$a_0$. Equation E12 still defines $a_0=M(U+1)$. Every main Pell
equation uses the mathematical parameter $A=a_0+4$, including
$G=1+(A+1)(f^2-1)$ in the signed auxiliary block and the modulus
$A-1$ in E20. The first Pell norm and index equations retain the forms
\begin{equation}\label{eq:base-four-first-pell}
 \tau(\tau+1)=Q(Q+1)k^2,\qquad k=R+1+hQ.
\end{equation}
The positive interval remains
\begin{equation}\label{eq:base-four-interval}
 c=kY_P+\eta,\qquad k=\eta+\zeta,
 \qquad\eta,\zeta>0,
 \qquad Y_P<\frac ck<Y_P+1.
\end{equation}
Replace E14 by
\begin{equation}\label{eq:base-four-congruence}
 d=W+c(A-4)+\gamma(8A-17).
\end{equation}
This is the source's $\chi$ exponent congruence for base four.
The second exponent congruence retains the form
\begin{equation}\label{eq:base-four-second-congruence}
 \mu=q+\kappa(A-B)+\rho(2AB-B^2-1).
\end{equation}

\subsection{Exact indices and the shifted approximation}

The pre-Pell part of Proposition~\ref{prop:positive-coefficient108}
uses only the encoding equations and gives
\[
 N=q^8,\qquad q>4b\ge8,\qquad
 3L\le B\le N\le R,\qquad b,q\le N.
\]
In particular $N,R\ge64$, $U,Y_P\ge4096$, $M\ge262144$,
$Q>R+1$, and $A>J>1$. These estimates precede both exponent
relations; the stronger $N\ge64$ will be needed for base four.

Since $h>0$, equation~\eqref{eq:base-four-first-pell} gives
$k\ge R+2$. The interval gives $c>Y_P(R+2)>J$.
The generic signed Pell argument of Proposition~\ref{lem:signed-pell}
therefore applies with parameter $A$ and gives
\begin{equation}\label{eq:base-four-main-index}
 c=\psi_A(J),\qquad d=\chi_A(J).
\end{equation}
The triangular norm gives the ordinary Pell equation for
$\widehat\tau=2\tau+1$ and $P=2Q+1$, so
$k=\psi_P(t)$ for some positive index $t$. The congruence
$\psi_P(t)\equiv t\pmod{2Q}$, reduced modulo $Q$, implies
\[
 t=R+1+vQ,\qquad v\ge0.
\]
If $v\ge1$, the elementary Pell growth bounds give
\[
 \frac ck\le
 \frac{(2A)^{2R}}{(2P-1)^{R+Q}}
 \le\left(\frac{2A}{2P-1}\right)^{2R}<\frac12.
\]
Here $Q\ge R$, and $2A/(2P-1)<1/2$ follows, for example, from
\[
 a_0<2UM,\qquad
 4A<8UM+16<4UM^2+1=2P-1.
\]
This contradicts~\eqref{eq:base-four-interval}. Consequently
\begin{equation}\label{eq:base-four-first-index}
 k=\psi_P(R+1),\qquad \widehat\tau=\chi_P(R+1).
\end{equation}
No exponent or rounding conclusion has been used to recover either
Pell index.

At these indices, the lower growth estimate gives
\begin{align}
 \frac ck
 &\ge\frac{(2A-1)^{2R}}{(2P)^R}\notag\\
 &=\xi\left(1+\frac7{2a_0}\right)^{2R}
          \left(1+\frac1{2Q}\right)^{-R}
 >\xi.\label{eq:base-four-lower}
\end{align}
Indeed $14Q>a_0$, so
$(1+7/(2a_0))^2>1+7/a_0>1+1/(2Q)$.
The upper growth estimate gives
\begin{align}
 \frac ck
 &\le\frac{(2A)^{2R}}{(2P-1)^R}
 <\xi\left(1+\frac4{a_0}\right)^{2R}\notag\\
 &<\xi\left(1+\frac{16R}{a_0}\right)
 =\xi\left(1+\frac{16}{Y_P(U+1)}\right).
 \label{eq:base-four-upper}
\end{align}
For the last step, binomial coefficients give
$(1+t)^m\le\sum_{j\ge0}(mt)^j=(1-mt)^{-1}<1+2mt$
when $0<mt<1/2$. Here $mt=8R/a_0=8/[Y_P(U+1)]<1/2$.

The lower bound and the interval imply $\xi<Y_P+1$.
Since $\xi>U^R$, integrality gives $Y_P\ge U^R$ and hence
\begin{equation}\label{eq:base-four-size}
 A>a_0=RY_P(U+1)>RU^R.
\end{equation}
This size bound is established before invoking an exponent lemma.
Also $\xi<Y_P+1<2Y_P$, so the upper bound yields
\begin{equation}\label{eq:base-four-error}
 0<\frac ck-\xi<\frac{32}{U+1}<\frac12.
\end{equation}

\subsection{The exponent relations and rounding}

Since $W=bw\le U$ and $N,R\ge64$, we have
\[
 4^{3J}=64\cdot64^{2R}
       \le64N^{2R}\le RU^R<A,
 \qquad W^3\le U^R<A.
\]
The $\chi$ version of source Lemma~2.22, explicitly stated after
that lemma, applies to~\eqref{eq:base-four-congruence} and
\eqref{eq:base-four-main-index}. It proves
\begin{equation}\label{eq:base-four-power}
 bw=4^J.
\end{equation}
Thus $b$ is a power of two, and
\[
 U=N^2w\ge Nbw=N4^J>4\cdot2^{2R}.
\]
E13, E19 and E20 give $0<\kappa<c$, its Pell norm, and
$\kappa\equiv L\pmod{A-1}$. Since $0<L<J<A$, source
Lemma~2.23 gives $\kappa=\psi_A(L)$ and $\mu=\chi_A(L)$.
Moreover
\[
 B^{3L}\le B^B\le N^N\le U^R<A,
 \qquad q^3\le N^N\le U^R<A.
\]
The same exponent lemma, applied now to
\eqref{eq:base-four-second-congruence}, proves $q=B^L$.

Let $F=\lfloor\xi\rfloor$. Source Lemma~2.24, or its direct
binomial expansion, gives
\begin{equation}\label{eq:base-four-fraction}
 0<\xi-F<\frac{2^{2R}}U<\frac14,
 \qquad F\equiv\binom{2R}{R}\pmod U.
\end{equation}
Combining the interval with~\eqref{eq:base-four-lower} and
\eqref{eq:base-four-error} gives
\[
 F<\xi<\frac ck<Y_P+1,
 \qquad Y_P<\frac ck<\xi+\frac12<F+\frac34.
\]
Both $F$ and $Y_P$ are integers. The first chain gives $F\le Y_P$;
the second gives $Y_P\le F$. Thus
\begin{equation}\label{eq:base-four-rounding}
 Y_P=F,\qquad N^2\mid\binom{2R}{R}.
\end{equation}
This argument requires no parity assertion. The power-of-two facts
for $b$ and $q$, the relation $q=B^L$, and the divisibility in
\eqref{eq:base-four-rounding} are precisely the conclusions needed
by the unchanged no-carry and modular-Sidon decoding proof.
They therefore imply membership in the represented r.e.\ set.

\subsection{Constructing positive witnesses}

For the converse, start from a solution of the represented system
and take its canonical modular-Sidon coding witnesses. Choose the
usual sufficiently large power of two $b$, put $B=Hb^2$,
$q=B^L$, $N=q^8$, and construct the unchanged packed $S,T,R$.
Their no-carry property gives $N^2\mid\binom{2R}{R}$.
These choices precede the Pell witnesses and supply $b\le R$
and $R\ge N\ge64$.

Put $J=2R+1$ and choose
\[
 w=4^J/b,\quad U=N^2w,\quad
 Y_P=\left\lfloor\frac{(U+1)^{2R}}{U^R}\right\rfloor,
 \quad s=Y_P/N^2.
\]
Writing $b=2^t$, the inequality $t\le b\le R<2J$ shows that
$w$ is a positive integer. The binomial expansion and central-binomial
divisibility make $s$ a positive integer. Define
$M=RY_P$, $a_0=M(U+1)$, $A=a_0+4$, $Q=UM^2$, $P=2Q+1$, and
\[
 c=\psi_A(J),\qquad d=\chi_A(J),\qquad k=\psi_P(R+1).
\]
The ratio estimates above hold at these chosen indices.
Here $\xi<Y_P+1$ is immediate from the definition of $Y_P$, so
\[
 Y_P<\xi<\frac ck<\xi+\frac12<Y_P+\frac34.
\]
Consequently
\[
 \eta=c-kY_P>0,\quad \zeta=k-\eta>0,\quad
 \tau=\frac{\chi_P(R+1)-1}{2},\quad
 h=\frac{k-R-1}{Q}
\]
are positive integers. Positivity and integrality of $\tau$ follow
from the oddness of the $\chi$ sequence for odd $P$; those of $h$
follow from the Pell congruence modulo $2Q$ and
$\psi_P(R+1)>R+1$. Thus the triangular norm, index equation,
positive interval, E12, and E15 all hold.

Choose positive $f,i$ with
$f^2-(A^2-1)(ic^2)^2=1$ using the same generic Pell divisibility
construction as in Proposition~\ref{lem:signed-pell}. Put
\[
 G=1+(A+1)(f^2-1),\quad I=\chi_G(J),\quad H_J=\psi_G(J),
 \qquad o=(I+d)/f,\quad j=(H_J-J)/c.
\]
Because $G\equiv-A\pmod f$, $G\equiv1\pmod c$, and $J$ is odd,
$o$ and $j$ are integers. They are positive since $I+d>0$ and
$\psi_G(J)>J$. These values satisfy the signed E17 with $of-d=I$.
Next choose
\[
 \kappa=\psi_A(L),\quad \mu=\chi_A(L),\quad
 \Delta=\frac{\kappa-L}{A-1},\quad \phi=c-\kappa.
\]
Since $J>L\ge2$, these are positive integers with the required
Pell norm, congruence and gap.

The exponent lemma supplies integer quotients $\gamma,\rho$ in
\eqref{eq:base-four-congruence} and
\eqref{eq:base-four-second-congruence}. To check their positivity,
use $\chi_A(t)=A\psi_A(t)-\psi_A(t-1)$ to obtain
\[
 d-(A-4)c>3c>W,
 \qquad \mu-(A-B)\kappa>(B-1)\kappa>q.
\]
For the last inequalities, $A>W^3,q^3$ and $J,L\ge2$ give
$c,\kappa\ge A$. The moduli are positive because $A>B\ge2$
and $A>4$. Hence $\gamma,\rho>0$.
All coding witnesses, including $\Omega$ and $\sigma$, remain
the canonical ones. This constructs every required positive witness.
In particular the construction re-chooses $w,s,a_0$ and their
dependent Pell coordinates; it does not preserve the predecessor's
base-two value of $w$.

\subsection{The complete arithmetic count}

\begin{theorem}\label{thm:best106}
For each r.e.\ set, the fixed modular-Sidon encoding gives a
membership-equivalent system in 34 positive unknowns and 22 equations
with a straight-line certificate of \textbf{106 arithmetic instructions}:
59 multiplications and 47 additions. Equality tests and supplied
parameters are free, as are fixed numerals.
\end{theorem}
\begin{proof}
The preceding arguments prove both directions over positive integers.
The arithmetic expressions are evaluated as
\[
\begin{gathered}
 A-4=a_0,\quad A-1=a_0+3,\quad A+1=(A-1)+2,\\
 A^2-1=(A-1)(A+1),\quad
 8A-17=8a_0+15,\quad A-B=a_0-(B-4).
\end{gathered}
\]
The previous subtraction computing the base-two coefficient $a-2$
disappears. The other expressions cost the same numbers of operations
as their predecessors; in particular $B-4$ is already computed by
the third mask, and the defining equation for $a_0$ retains its
existing factorization. No register for $A$ itself is required.

The checker \file{round14\_1980\_base\_four\_certificate.py}
verifies all source residuals, both triangular corrections, every
shifted-register identity, and the full primitive instruction list.
It counts 59 multiplications and 47 additions. The positive input
and equality counts are unchanged. Its generated schedule records
the exponent $L=2^{32}$.
\end{proof}
